\section{Introduction}
\label{sec:intro}

The rapid proliferation of specialized deep neural networks has ushered in a new era of machine learning, where foundation models are routinely fine-tuned on diverse, downstream tasks to create specialized expert models. While these expert models exhibit exceptional performance within their targeted domains, deploying them in resource-constrained environments poses severe challenges due to the prohibitive memory and computational overhead of hosting multiple distinct networks. Consequently, consolidating these task-specific capabilities into a single, multi-task model has emerged as a critical research frontier. Joint multi-task retraining~\cite{caruana1997multitask} offers a potential solution but is often practically unviable due to high computational costs, data privacy constraints, and the risk of catastrophic forgetting of previous knowledge.

To bypass these limitations, post-hoc model merging has emerged as a highly popular and resource-efficient paradigm. Pioneered by works such as Task Arithmetic~\cite{ilharco2022editing} and Model Soups~\cite{wortsman2022model}, post-hoc model merging constructs a multi-task network by linearly combining the weights of pre-trained expert models. However, simple uniform weight interpolation often suffers from severe representation clashes and inter-task interference, degrading performance across all tasks. To resolve these conflicts, recent state-of-the-art methods have introduced parameterized merging, which seeks to optimize layer-wise or parameter-wise scaling coefficients. Methods such as AdaMerging~\cite{adamerging} and RegCalMerge~\cite{regcalmerge} optimize these coefficients using test-time data streams or small offline validation sets.

Despite their empirical success, we argue that existing parameterized merging frameworks are fundamentally constrained by a shared, unexamined assumption: they operate and optimize strictly within the physical, spatial depth coordinate space. Under this spatial paradigm, layer-wise coefficients are treated as independent, unconstrained variables. We hypothesize that this representation is highly redundant and poorly conditioned. Because adjacent neural layers exhibit strong functional coupling, the optimal merging coefficients should follow a smooth, continuous trajectory across network depth. Unconstrained spatial optimization ignore this underlying structure, allowing the coefficients to oscillate wildly. Under data-scarce (few-shot) regimes, these high-frequency spatial oscillations behave as high-variance optimization noise, leading to catastrophic overfitting to the local validation stream—a phenomenon we term the \emph{Overfitting-Optimizer Paradox}.

In this work, we propose to rethink this spatial coordinate paradigm. We pose a visionary question: \emph{What if layer-wise parameter sensitivity is best modeled as a multi-scale spectral distribution rather than a spatial sequence?} Since a deep neural network is organized as a discrete 1D sequence of layers, the layer-wise merging coefficient profile is mathematically equivalent to a discrete 1D spatial signal. Drawing inspiration from digital signal processing (DSP), we can map this spatial signal into the frequency domain using the Discrete Cosine Transform (DCT-II). Under this spectral transformation, the low-frequency coefficients represent the macro-structural, task-sharing pathways that span across the network's depth, while high-frequency coefficients represent micro-structural, task-specific local variations.

By transitioning parameter consolidation into the frequency domain, we present \textbf{SpectralMerge}, a novel, mathematically elegant model merging framework. SpectralMerge allows us to directly control and regularize the spectral composition of merging coefficients. We formulate and evaluate two distinct frequency-domain constraints:
\begin{enumerate}
    \item \textbf{SpectralMerge-LP (Low-Pass Hard Cutoff):} We restrict the trainable parameters to the first $F$ low-frequency DCT components, forcing all higher-frequency coefficients to zero. This acts as an analytical low-pass filter that mathematically prevents the optimizer from fitting high-frequency spatial noise.
    \item \textbf{SpectralMerge-Reg (Soft Spectral Regularization):} We optimize all frequency components but add a quadratic **Spectral Decay Penalty** to the loss function, which penalizes high-frequency oscillations progressively.
\end{enumerate}
Importantly, by leveraging the orthogonal trigonometric basis functions of the DCT, SpectralMerge enjoys perfect numerical conditioning. This stands in sharp contrast to spatial smoothing techniques such as PolyMerge~\cite{polymerge}, which uses low-degree polynomials that suffer from severe collinearity and boundary-condition instability.

Through comprehensive evaluations across 30 random seeds, we demonstrate the exceptional performance and robustness of SpectralMerge. To ensure a controlled and mathematically rigorous analysis of the optimization and generalization dynamics, our evaluations leverage a continuous multi-task weight-merging simulation landscape calibrated directly on Vision Transformer (ViT-B/32) empirical sensitivity statistics. On standard clean streams, SpectralMerge-LP ($F=3$) and SpectralMerge-Reg achieve state-of-the-art simulated accuracies of $85.32\%$ and $85.17\%$, outperforming both online adaptation baselines and spatial polynomial models. Under extreme data scarcity ($M \le 10$), SpectralMerge completely resolves the Overfitting-Optimizer Paradox, outperforming unconstrained spatial search by over $2.6\%$ on the simulation benchmark. Furthermore, we demonstrate that our frequency-domain parameterization is immune to validation target selection bias and adversarial stream noise, making it highly suitable for real-world deployment.

In summary, our key contributions are:
\begin{itemize}
    \item We challenge the traditional spatial coordinate paradigm of model merging and propose a novel frequency-domain parameterization via the Discrete Cosine Transform (DCT-II).
    \item We introduce SpectralMerge-LP and SpectralMerge-Reg, two mathematically elegant frequency-domain formulations that provide robust structural regularization and perfect numerical conditioning.
    \item We empirically demonstrate that SpectralMerge outpeforms all competitive baselines on standard clean streams and across multiple adversarial conditions (extreme label shift, bursty streams, and batch noise).
    \item We systematically refute the Overfitting-Optimizer Paradox under extreme sample complexity, and prove that frequency-domain model merging is highly resilient to validation target selection bias.
\end{itemize}
